\documentclass[conference]{IEEEtran}
\IEEEoverridecommandlockouts
% The preceding line is only needed to identify funding in the first footnote. If that is unneeded, please comment it out.
\usepackage{cite}
\usepackage{amsmath,amssymb,amsfonts}
\usepackage{algorithmic}
\usepackage{graphicx}
\usepackage{textcomp}
\usepackage{xcolor}
\usepackage{caption}
\usepackage{float}
\usepackage{hyperref}
\def\BibTeX{{\rm B\kern-.05em{\sc i\kern-.025em b}\kern-.08em
    T\kern-.1667em\lower.7ex\hbox{E}\kern-.125emX}}

\begin{document}

\title{Psychologist AI Agent\\
{\footnotesize \textsuperscript{*}EECS E6895 Advanced Big Data and AI Midterm Project}
}

\author{
    \IEEEauthorblockN{Changyuan Yu}
    \IEEEauthorblockA{\textit{Columbia University} \\
    cy2812@columbia.edu}
    \and
    \IEEEauthorblockN{Zhiliang Wang}
    \IEEEauthorblockA{\textit{Columbia University} \\
    zw3162@columbia.edu}
    \and
    \IEEEauthorblockN{Yihang (Kevin) Sun}
    \IEEEauthorblockA{\textit{Columbia University} \\
    ys3978@columbia.edu} 
}

\maketitle

\begin{center}
    \textbf{\large Abstract}
\end{center}
\noindent The rapid integration of Large Language Models (LLMs) into mental health services has highlighted critical tensions between conversational capability, clinical reliability, and data privacy. Generic cloud-based models often falter under the risk of personal information leakage and ``hallucinations'' that can compromise patient safety. This paper introduces a professional, privacy-preserving AI Psychologist Agent characterized by a hybrid ``Cloud-Analysis, Local-Generation'' architecture. By leveraging the reasoning depth of Deepseek-V3 to guide a locally deployed, DPO-optimized Llama-3.1-8B model, the system ensures that sensitive user data remains on-device while maintaining high-tier therapeutic intelligence. Our methodology integrates an offline knowledge distillation pipeline with an online inference framework featuring a multi-layered safety gateway. Specifically, the agent employs a Retrieval-Augmented Generation (RAG) system grounded in CBT and DBT frameworks to minimize clinical errors, alongside a reversible PII redaction mechanism to shield identity during cloud-assisted auditing. Experimental results demonstrate that the proposed ``Only-Escalate'' risk policy and semantic safety filters achieve high accuracy in intercepting self-harm and medical non-compliance triggers.

\begin{IEEEkeywords}
AI Psychologist, Privacy-Preserving, Knowledge Distillation, RAG, Large Language Models
\end{IEEEkeywords}

\section{Introduction}

The rapid digitalization of mental health services has created an unprecedented demand for accessible AI-driven psychological support. As professional therapy remains costly and geographically constrained, Large Language Models (LLMs) offer a promising alternative for providing 24/7 empathetic intervention. However, the transition from generic conversational AI to a specialized clinical psychologist agent faces three formidable hurdles.

First, the \textit{privacy-utility trade-off} remains a primary concern. Mental health dialogues are inherently dense with Personal Identifiable Information (PII). Traditional cloud-based LLM services require transmitting this sensitive data to external servers, raising significant risks of data exfiltration and violating the "privacy-first" principle of psychological counseling. Second, \textit{clinical reliability} is often compromised by "hallucinations," where models generate factually incorrect or therapeutically unsound advice. In a high-stakes clinical context, such errors are not merely technical glitches but potential safety hazards. Third, \textit{safety boundary management} in existing models is often either too restrictive (refusing to help) or too lax (failing to detect crisis triggers like self-harm).

To bridge these gaps, this project proposes a professional, privacy-preserving Psychologist AI Agent. Our core innovation lies in a hybrid ``Cloud-Analysis, Local-Generation'' architecture. This design philosophy delegates complex clinical reasoning to a powerful cloud model (Deepseek-V3) using redacted, non-sensitive text, while the actual empathetic response is generated by a locally deployed, Direct Preference Optimization (DPO) tuned Llama-3.1-8B model. 

The primary contributions of this work are summarized as follows:
\begin{itemize}
    \item We developed a dual-pipeline framework that integrates offline knowledge distillation from Deepseek-V3 to Llama-3.1-8B, enhancing the local model's clinical professionality through DPO.
    \item We implemented a multi-layered safety and privacy infrastructure, including a semantic safety gateway, a reversible PII redaction module, and an ``Only-Escalate'' risk auditing policy.
    \item We integrated a Retrieval-Augmented Generation (RAG) system grounded in professional psychological frameworks such as CBT, DBT, and WHO guidelines to minimize hallucinations and ensure evidence-based responses.
    \item Experimental evaluations demonstrate that our agent achieves high-precision risk interception and maintains clinical accuracy on consumer-grade hardware (8GB VRAM), offering a scalable solution for private mental health support.
\end{itemize}

The remainder of this paper is organized as follows. Section II details the proposed system architecture and its core modules. Section III describes the implementation of our offline and online pipelines. Section IV presents the evaluation results and demo scenarios, followed by future work and conclusions in Section V and VI.

\section{Methodology}
The architecture of the Psychologist AI Agent is founded on a bifurcated design, meticulously engineered to balance high-level clinical reasoning with uncompromising data privacy. This framework is divided into two primary stages: an offline knowledge distillation and preference alignment phase for model optimization, and a multi-step online inference pipeline for secure, real-time interaction.

\subsection{Offline Optimization: Knowledge Distillation and DPO}
The development of a specialized local model begins with a rigorous data refinement and distillation process, aimed at Narrowing the gap between generic LLM capabilities and professional psychological expertise. Our initial dataset comprised 2,775 counselor-patient dialogues, which underwent a strict cleaning and deduplication protocol to ensure clinical relevance. This resulted in a core set of 863 high-quality samples, representing a 31\% retention rate. To further address the scarcity of data in specialized clinical topics, we leveraged Deepseek-V3 to synthesize an additional 194 scenarios, creating a robust training corpus of 1,057 dialogues. 
\begin{figure}[htbp]
    \centering
    \includegraphics[width=\linewidth]{Pipeline1.png} 
    \caption{Offline Pipeline}
    \label{fig:pipeline1}
\end{figure}
Central to our model’s professional alignment is the implementation of Direct Preference Optimization (DPO). Unlike traditional supervised fine-tuning, the DPO pipeline utilizes Deepseek-V3 as a "Professional Judge" to construct a preference dataset. For each clinical prompt, the system pairs a sophisticated, empathetic response (the chosen response) against a baseline, often overly generic or clinically shallow response from a standard Llama model (the rejected response). This alignment process, executed via the Unsloth framework using QLoRA, allows the locally deployed Llama-3.1-8B model to internalize the nuanced linguistic styles and decision-making logic of an expert therapist. The final model is exported into a 4-bit GGUF format, specifically optimized for high-speed inference on consumer-grade 8GB VRAM hardware.

\subsection{Online Inference: The Privacy-Preserving Pipeline}
The online inference engine is structured as an 8-step orchestrated sequence designed to prevent the exfiltration of sensitive information while maximizing therapeutic utility. The process initiates with a Semantic Safety Gateway that utilizes BGE-small embeddings to compare user inputs against 217 malicious or non-compliant patterns. By setting a calibrated semantic threshold between 0.75 and 0.85, the system identifies potential risks before they are even processed for reasoning. Concurrently, a Personal Identifiable Information (PII) redaction module, integrating Microsoft Presidio with customized regular expressions, masks sensitive entities such as names and locations. This ensures that only de-identified, redacted text is ever transmitted to the cloud-based supervisor.
\begin{figure}[htbp]
    \centering
    \includegraphics[width=\linewidth]{Pipeline2.png} 
    \caption{Online Pipeline}
    \label{fig:pipeline2}
\end{figure}


A critical innovation in our pipeline is the dual-prompt supervision mechanism. Once the user's intent is de-identified and contextualized through a FAISS-based Retrieval-Augmented Generation (RAG) search across CBT and DBT manuals, the system sends a diagnostic request to the cloud-based Deepseek-V3. This "Cloud Supervisor" performs a high-level clinical analysis and safety audit, returning structured guidance without ever seeing the user's actual identity. Finally, the local agent receives this professional blueprint and generates the final response using the original, unmasked text stored in local memory. This "Only-Escalate" policy ensures that if the cloud audit detects high-risk triggers, such as self-harm ideation, the local generation is overridden by an immediate, hard-coded crisis intervention protocol. This hybrid synergy allows the agent to maintain a deep therapeutic alliance while operating within a strictly guarded privacy boundary.

\section{Evaluation and Results}
To rigorously validate the efficacy of the Psychologist AI Agent, we conducted a multidimensional evaluation focusing on four critical aspects: semantic safety, retrieval precision, human-centric clinical quality, and real-time inference efficiency.

\subsection{Safety Gateway Benchmarking}
The primary objective of the safety evaluation is to assess the system's resilience against high-risk inputs that frequently challenge generic LLMs, such as self-harm ideation, medical non-compliance, and violent triggers. We utilized a curated benchmark containing 217 malicious attack patterns and evaluated the interception rate of our BGE-small embedding-based gateway. The methodology involved testing various semantic similarity thresholds (ranging from 0.75 to 0.85) to find the optimal balance between security and conversational fluidity. 

The results demonstrated a 100\% interception rate for explicit self-harm and medical advice queries at the 0.85 threshold. Analysis of the data suggests that while a higher threshold ensures a "zero-tolerance" policy for critical risks, the integration of an "Only-Escalate" secondary audit by the cloud supervisor provides a necessary safety net for ambiguous cases. This multi-layered defense ensures that the agent remains within clinical ethical boundaries without prematurely terminating benign user interactions.
\begin{figure}[htbp]
    \centering
    \includegraphics[width=\linewidth]{Safety.png} 
    \caption{Safety Gateway}
    \label{fig:Safety Gateway}
\end{figure}

\subsection{RAG Pipeline Precision}
Given the high stakes of mental health support, eliminating "hallucinations" is paramount. We evaluated the Retrieval-Augmented Generation (RAG) module to determine its ability to ground agent responses in professional clinical frameworks (CBT/DBT). The evaluation involved 215 test cases mapped against 179 knowledge chunks from WHO and professional therapeutic manuals. We measured performance using Top-k hit rates and Mean Reciprocal Rank (MRR). 

The RAG system achieved a 100\% Top-5 hit rate and an MRR of 0.875, indicating that the FAISS-based vector search consistently identifies the most relevant clinical guidelines. The analysis confirms that this retrieval precision directly translates to higher factual accuracy in the local model’s generation, as the injected context effectively constrains the LLM’s tendency to improvise clinically unverified advice.
\begin{figure}[htbp]
    \centering
    \includegraphics[width=\linewidth]{RAG Top5.png} 
    \caption{RAG Precision}
    \label{fig:RAG Precision}
\end{figure}

\subsection{Agent Quality Assessment}
Beyond quantitative metrics, the therapeutic alliance depends on empathy and clinical nuance. We implemented a "Human-in-the-Loop" evaluation where Deepseek-V3 acted as an expert auditor to score the agent's dialogues on a scale of 1--10 across three dimensions: Empathy, Clinical Accuracy, and Safety Compliance. This method was chosen to simulate the subjective judgment of a professional counselor.

The DPO-optimized Llama-3.1 model outperformed the baseline model significantly, achieving an average Empathy score of 9.2 and a Clinical Accuracy score of 8.8. Our analysis reveals that the Direct Preference Optimization (DPO) phase successfully filtered out the "robotic" and overly-cautious tone of the base model, replacing it with validated, empathetic responses. This improvement demonstrates that our distillation pipeline effectively transfers specialized therapeutic intuition from a large-scale teacher model to a compact student model.
\begin{figure}[htbp]
    \centering
    \includegraphics[width=\linewidth]{Agent2.png} 
    \caption{Response Score}
    \label{fig:Response Score}
\end{figure}

\subsection{Inference and Latency Performance}
Finally, we evaluated the system's operational efficiency to ensure its viability for local deployment on consumer-grade hardware (8GB VRAM). The evaluation focused on end-to-end latency and VRAM consumption across the full 8-step inference pipeline. Performance was measured during peak concurrent processing to simulate real-world usage.

The experimental results showed that the safety gateway maintains an average latency of only 86.99ms, with the entire pipeline delivering responses with minimal overhead. By utilizing 4-bit GGUF quantization and full GPU offloading, the system comfortably fits within the 8GB VRAM limit while maintaining high throughput. This proves that the hybrid architecture successfully mitigates the computational burden of complex auditing, making professional-grade AI counseling accessible on standard personal computers without sacrificing privacy or speed.
\begin{figure}[htbp]
    \centering
    \includegraphics[width=\linewidth]{Agent1.png} 
    \caption{Latency}
    \label{fig:Latency}
\end{figure}

\section{Demo}
We demonstrate the system through a multi-step interaction highlighting four core capabilities: context understanding, therapeutic reasoning, crisis detection, and memory-based context tracking.

\textbf{Initial Interaction}:
The user first introduces themselves and describes stress related to upcoming exams. The system responds with an empathetic message while the pipeline performs retrieval, supervisor analysis, and local generation to identify academic stress as the primary concern.
\begin{figure}[htbp]
    \centering
    \includegraphics[width=0.23\linewidth]{6895-1.jpg}
    \includegraphics[width=0.23\linewidth]{6895-2.jpg}
    \caption{Initial Interaction}
    \label{fig:Initial Interaction}
\end{figure}

\textbf{Therapeutic Reasoning}:
When the user expresses catastrophic thinking about failing the exam, the system identifies this cognitive distortion and applies CBT-style reasoning. The generated response validates the user’s feelings while encouraging cognitive reframing and healthier coping strategies.
\begin{figure}[htbp]
    \centering
    \includegraphics[width=0.23\linewidth]{6895-3.jpg}
    \includegraphics[width=0.23\linewidth]{6895-4.jpg}
    \caption{Therapsutic Reasoning}
    \label{fig:Therapautic Reasoning}
\end{figure}

\textbf{Crisis Detection}:
When the user expresses suicidal ideation, the safety gateway and supervisor classify the message as high risk. The risk audit activates the crisis-response protocol and the system provides emergency support resources such as the 988 hotline.
\begin{figure}[htbp]
    \centering
    \includegraphics[width=0.23\linewidth]{6895-5.jpg}
    \includegraphics[width=0.23\linewidth]{6895-6.jpg}
    \caption{Crisis Detection}
    \label{fig:Crisis Detection}
\end{figure}

\textbf{Memory and Context Recall}:
Finally, the user asks whether the system remembers the previously mentioned exam. The agent retrieves this information from the user profile memory module, demonstrating its ability to maintain longitudinal conversational context.
\begin{figure}[htbp]
    \centering
    \includegraphics[width=0.23\linewidth]{6895_7.jpg}
    \includegraphics[width=0.23\linewidth]{6895_8.jpg}
    \caption{Long-term Memory}
    \label{fig:Long-term Memory}
\end{figure}

Our Github repository link: 
\url{https://github.com/yuchangyuan1/Psychologist_Agent}

Our demo video link: 
\url{https://drive.google.com/file/d/1xyuwFvmP21YAC60IsNfKYa-sQwqbUf6A/view?usp=sharing}

\section{Conclusion and Future Work}
This paper presented a professional, privacy-preserving AI Psychologist Agent that bridges the gap between high-level clinical reasoning and local data security. By implementing a hybrid architecture that delegates analytical auditing to the cloud while keeping empathetic generation on-device, we have created a scalable framework for sensitive healthcare applications. The integration of DPO-optimized knowledge distillation and RAG-based grounding ensures that the agent is not only safe and private but also clinically informed. Our benchmarks confirm that the system is capable of high-precision risk detection and empathetic engagement on consumer-grade hardware, making professional mental health support more accessible.

Looking ahead, we intend to refine the agent's longitudinal memory capabilities by implementing a "living document" UserProfile module, which will allow for cross-session progress tracking. Additionally, we plan to explore more advanced alignment techniques, such as Generalized EM-based Preference Alignment (Gepa), to further calibrate the model's responses with expert-level therapeutic intuition. Finally, incorporating multi-modal emotional analysis, such as voice tone and facial micro-expression recognition, remains a key objective for achieving a truly holistic digital therapeutic experience.

\end{document}